\section{Readiness Reporting and Operations}
\label{sec:operations}
Afar-prgenv turns monthly AFAR drops into auditable support artifacts. The
repro harness emits per-profile logs, and the P/S/F outcomes are summarized in
readiness reports that track drop stability across time. This operational view
is important for user support because it distinguishes regressions introduced
by new drops from persistent gaps in the PE ecosystem.

Readiness reports capture the AFAR drop ID, PE version, ROCm version, and a
summary of PASS/SKIP/FAIL counts, plus the explicit reason for each SKIP. This
keeps dependency gaps visible (for example, missing ROCm 6 runtime libraries
or incomplete `pkg-config` metadata) while allowing early adopters to
experiment with drops that otherwise meet functional readiness criteria.

\subsection{Readiness Labels}
\label{sec:labels}
Each repro outcome is classified as PASS, SKIP, or FAIL. PASS indicates that
the workflow compiles and links with the PE wrappers and required libraries.
SKIP indicates a known precondition is missing (for example, a runtime
dependency or missing metadata) and is treated as a gate rather than a
regression. FAIL indicates a functional break in the drop or wrapper layer and
blocks progression to extended validation.

\begin{table}[t]
  \centering
  \small
  \setlength{\tabcolsep}{4pt}
  \caption{Readiness labels and operational response.}
  \label{tab:labels}
  \begin{tabular}{l p{0.62\columnwidth}}
    \hline
    Label & Operational response \\
    \hline
    PASS & Proceed to extended validation and user-facing evaluation. \\
    SKIP & Document precondition, track dependency, rerun when resolved. \\
    FAIL & Block release; isolate regression and update wrapper or metadata. \\
    \hline
  \end{tabular}
\end{table}

\subsection{Readiness Report Contents}
\label{sec:report}
Each drop evaluation produces a short readiness report that captures the
module stack, test coverage, and gating outcomes used in the decision. The
report format is intentionally lightweight so it can be generated monthly and
shared with operations staff.

\begin{table}[t]
  \centering
  \small
  \setlength{\tabcolsep}{4pt}
  \caption{Core fields in a readiness report.}
  \label{tab:report-fields}
  \begin{tabular}{l p{0.62\columnwidth}}
    \hline
    Field & Purpose \\
    \hline
    AFAR drop version & Track the compiler/toolchain drop under evaluation. \\
    PE module stack & Record \texttt{cpe}, \texttt{PrgEnv}, MPI, and targets. \\
    ROCm version(s) & Capture primary ROCm and legacy dependencies. \\
    Repro coverage & List profiles and module-order permutations executed. \\
    P/S/F summary & Summarize readiness outcomes plus explicit SKIP causes. \\
    LDD audit status & Note dependency allowlist results for binaries. \\
    \hline
  \end{tabular}
\end{table}

\subsection{Cadence and Regression Tracking}
\label{sec:cadence}
AFAR drops arrive on a rapid cadence, so readiness reporting emphasizes
repeatability. The generator logs, repro outputs, and module versions are
captured together so that support teams can compare a new drop against the last
known-good baseline. This makes it easier to attribute regressions to a change
in compiler behavior versus a change in PE library metadata.

Because the reports are structured, we can trend the pass rate across drops
and flag persistent SKIP causes that should be addressed by PE library updates
or wrapper fixes. Over time, this reduces the time to certify each monthly
drop and clarifies which regressions are actionable by the support team.

\subsection{From Repro Gating to User Rollout}
\label{sec:rollout}
Once the repro suite passes and known SKIP conditions are documented, support
teams can stage the drop for early adopters. Additional validation suites
(SOLLVE V\&V and the OLCF Fortran applications) are used to confirm readiness
for broader user adoption, but the repro gate remains the minimum requirement
for publishing a drop in the site module tree.
